% =====================================================================
% PROJETO DE ENGENHARIA DE SOFTWARE — AULA 4 (23/11)
% Escolher a solução e organizar os dados
% Caso 15 · Reserva de laboratórios da faculdade
% O diagrama fica em figuras/ e é incluído com \input.
% =====================================================================

\documentclass[a4paper,12pt]{article}

\usepackage[utf8]{inputenc}
\usepackage[T1]{fontenc}
\usepackage[brazilian]{babel}
\usepackage{times}
\usepackage[a4paper,top=3cm,left=3cm,right=3cm,bottom=2.5cm]{geometry}
\usepackage{xltabular}
\usepackage{array}
\usepackage{amssymb}
\usepackage{titlesec}
\usepackage{enumitem}
\usepackage{tikz}
\usetikzlibrary{arrows.meta,positioning,shapes.multipart}
\usepackage{caption}

\setlength{\parindent}{0pt}
\setlength{\parskip}{6pt}
\renewcommand{\arraystretch}{1.3}
\setlist{itemsep=0pt,topsep=2pt}
\captionsetup{font=small,labelfont=bf,labelsep=period}

\titleformat{\section}{\fontsize{13}{15}\bfseries}{\thesection.}{1em}{}
\titleformat{\subsection}{\fontsize{12}{14}\bfseries}{\thesubsection.}{1em}{}

\newcolumntype{Y}{>{\raggedright\arraybackslash}X}

\begin{document}
\pagestyle{plain}

\begin{center}
  {\small\bfseries PROJETO DE ENGENHARIA DE SOFTWARE}\\[4pt]
  {\fontsize{16}{20}\selectfont\bfseries Aula 4 · 23/11 · Escolher a solução e organizar os dados}\\[6pt]
  {\large Caso 15 · Reserva de laboratórios da faculdade}
\end{center}

\vspace{6pt}

\begin{xltabular}{\linewidth}{|l|Y|}
  \hline
  \textbf{Professor} & Sérgio Santana \\ \hline
  \textbf{Disciplina} & Laboratório de Engenharia de Software \\ \hline
  \textbf{Integrante} & Wilia Francisco da Silva \\ \hline
  \textbf{Regra para testar} & Reservas de 1 a 4 horas, para até 40 alunos \\ \hline
\end{xltabular}

% ---------------------------------------------------------------------
\section{Comparação de duas formas de construir o sistema}

Foram comparadas duas opções para a versão 2.0: um \emph{aplicativo
instalado} no celular (Android e iPhone) e um \emph{site que abre no
celular}, isto é, um site responsivo que se adapta à tela do celular e
do computador.

\begin{xltabular}{\linewidth}{|>{\raggedright\arraybackslash}p{2.9cm}|Y|Y|}
  \hline
  \textbf{Critério} & \textbf{Aplicativo instalado} & \textbf{Site que abre no celular} \\ \hline
  \endhead
  Custo & Mais alto: é preciso fazer e manter versões para Android e iPhone e pagar as contas das lojas (Google Play: taxa única de US\$ 25; Apple: US\$ 99 por ano). & Mais baixo: um único código serve para celular e computador; o custo principal é a hospedagem. \\ \hline
  Facilidade para o usuário & O professor precisa encontrar, baixar e instalar o aplicativo antes da primeira reserva; ocupa espaço para um uso de poucas vezes por semana. & Abre por um link, sem instalar nada. Funciona no celular e no computador da sala de aula. \\ \hline
  Manutenção & Cada correção passa pela revisão das lojas e só chega a quem atualizar o aplicativo; podem existir várias versões em uso ao mesmo tempo. & A correção é feita no servidor e vale na hora para todos os usuários; existe uma só versão em uso. \\ \hline
  Uso pela coordenação & A coordenação trabalha no computador; seria preciso fazer também uma versão para computador. & O mesmo site atende professores no celular e a coordenação no computador. \\ \hline
  Avisos ao professor & Pode enviar notificações no próprio celular. & Envia os avisos por e-mail (RF10), o que já atende ao caso. \\ \hline
\end{xltabular}

% ---------------------------------------------------------------------
\section{Solução escolhida e justificativa}

A solução escolhida é o \textbf{site que abre no celular}.

O problema do caso é que a coordenação recebia os pedidos de laboratório
por e-mail e dois professores chegavam para o mesmo horário. Para
resolver isso, todos precisam consultar e alterar \emph{a mesma agenda,
ao mesmo tempo}: quando um professor reserva o Lab 101, o horário deve
aparecer ocupado na hora para todos os outros. Um site ligado a um único
servidor garante isso de forma simples, para quem está no celular e para
quem está no computador.

Além disso, o professor reserva laboratório poucas vezes por semana.
Exigir que ele instale um aplicativo cria uma barreira que pode fazê-lo
voltar a pedir por e-mail, que é exatamente o risco R04 da Aula 2 e o
contrário do RNF01 (fazer uma reserva em menos de 2 minutos, sem
treinamento). O site também tem menor custo e manutenção mais simples,
e a versão 1.0 já funciona no navegador, o que facilita a migração dos
dados (risco R01). A única vantagem clara do aplicativo, as notificações
no celular, é coberta pelos avisos por e-mail previstos no RF10.

% ---------------------------------------------------------------------
\section{Informações que o sistema guarda}

As informações foram levantadas a partir dos campos do formulário
\emph{Nova reserva} e das colunas da tabela \emph{Registros} da versão
1.0, e agrupadas por assunto. A Figura~\ref{fig:dados} mostra os grupos
ligados por setas; o sinal \textbf{+} marca o que precisa ser
acrescentado na versão 2.0.

\begin{center}
  % Informações que o sistema guarda, agrupadas por assunto — Caso 15
% "+" marca o que precisa ser acrescentado na versão 2.0;
% grupos com borda tracejada são novos na versão 2.0.
\begin{tikzpicture}[
  grupo/.style={rectangle split, rectangle split parts=2, draw, thick, rounded corners=2pt,
                text width=3.55cm, align=left, font=\small, inner sep=4pt},
  novo/.style={dashed},
  liga/.style={-{Stealth[length=2.6mm]}, thick},
  rot/.style={font=\footnotesize, fill=white, inner sep=1.5pt, align=center}
]
  \node[grupo, fill=orange!8] (res) {%
    \textbf{Reserva}
    \nodepart{second}
    código\\
    professor {\footnotesize(na v2.0: usuário)}\\
    disciplina\\
    laboratório\\
    data\\
    início\\
    duração (horas)\\
    término (calculado)\\
    número de alunos\\
    situação (Reservado ou Cancelada)\\
    criada em\\
    \textbf{+} alterada em\\
    \textbf{+} cancelada em};

  \node[grupo, novo, fill=blue!6, left=1.35cm of res.north west, anchor=north east] (usu) {%
    \textbf{Usuário} {\footnotesize(novo na v2.0)}
    \nodepart{second}
    \textbf{+} código\\
    \textbf{+} nome\\
    \textbf{+} e-mail\\
    \textbf{+} senha (criptografada)\\
    \textbf{+} perfil (professor ou coordenação)};

  \node[grupo, fill=green!7, right=1.35cm of res.north east, anchor=north west] (lab) {%
    \textbf{Laboratório}
    \nodepart{second}
    nome (ex.: Lab 101)\\
    capacidade (30 ou 40 lugares)};

  \node[grupo, novo, fill=red!5, below=1.4cm of lab.south, anchor=north] (avi) {%
    \textbf{Aviso} {\footnotesize(novo na v2.0)}
    \nodepart{second}
    \textbf{+} código\\
    \textbf{+} reserva\\
    \textbf{+} e-mail de destino\\
    \textbf{+} tipo (confirmação ou cancelamento)\\
    \textbf{+} enviado em};

  \draw[liga] (usu.east |- res.text split) -- node[rot, above=2pt]{faz\\1:N} (res.west |- res.text split);
  \draw[liga] (res.east |- lab.text split) -- node[rot, above=2pt]{usa\\N:1} (lab.west |- lab.text split);
  \draw[liga] (res.east |- avi.text split) -- node[rot, above=2pt]{gera\\1:N} (avi.west |- avi.text split);

  \node[font=\footnotesize, anchor=north west, align=left] at (usu.south west |- res.south) {%
    \textbf{+} = acrescentar na versão 2.0\\
    borda tracejada = grupo novo na versão 2.0\\
    1:N = um para muitos; N:1 = muitos para um};
\end{tikzpicture}

  \captionof{figure}{Informações que o sistema guarda, agrupadas por assunto.}
  \label{fig:dados}
\end{center}

\begin{xltabular}{\linewidth}{|l|>{\raggedright\arraybackslash}p{3.4cm}|Y|Y|}
  \hline
  \textbf{Grupo} & \textbf{Informação} & \textbf{Onde aparece na versão 1.0} & \textbf{Versão 2.0} \\ \hline
  \endhead
  Reserva & Professor & Campo do formulário e coluna \emph{Professor} & Passa a vir do usuário que fez login (RF09). \\ \hline
  Reserva & Disciplina & Campo do formulário e coluna \emph{Disciplina} & Mantém. \\ \hline
  Reserva & Laboratório & Campo (lista) e coluna \emph{Laboratório} & Mantém, ligado ao grupo Laboratório. \\ \hline
  Reserva & Data & Campo do formulário e coluna \emph{Data} & Mantém. \\ \hline
  Reserva & Início & Campo (lista) e coluna \emph{Horário} & Mantém. \\ \hline
  Reserva & Duração em horas & Campo do formulário & Mantém. \\ \hline
  Reserva & Término & Calculado (início + duração) e mostrado na coluna \emph{Horário} & Mantém, calculado. \\ \hline
  Reserva & Número de alunos & Campo do formulário e coluna \emph{Alunos} & Mantém. \\ \hline
  Reserva & Situação & Coluna \emph{Situação} (Reservado ou Cancelada) & Mantém. \\ \hline
  Reserva & Data de criação & Guardada pelo sistema, sem aparecer na tela & Mantém. \\ \hline
  Reserva & Datas de alteração e de cancelamento & Não existem & \textbf{Acrescentar} (RF08 e RF03). \\ \hline
  Laboratório & Nome e capacidade & Fixos no sistema: Lab 202 com 30 lugares e os demais com 40 & \textbf{Acrescentar} um cadastro próprio, para mudar a capacidade sem alterar o código. \\ \hline
  Usuário & Nome, e-mail, senha e perfil & Não existem & \textbf{Acrescentar} (RF09 e RNF03). \\ \hline
  Aviso & Reserva, e-mail, tipo e data de envio & Não existem & \textbf{Acrescentar} (RF10). \\ \hline
\end{xltabular}

% ---------------------------------------------------------------------
\section{Artigo: fundamentação teórica}

\textbf{Tema 6 · Teste de usabilidade em aplicativos}

\subsection*{Usabilidade}

Barbosa e Silva (2010) apresentam a usabilidade como um dos principais
critérios de qualidade de uso de um sistema interativo. Com base na
norma ISO 9241-11, os autores a descrevem como o grau em que um produto
pode ser usado por usuários específicos para alcançar objetivos
específicos com eficácia, eficiência e satisfação, em um contexto de uso
determinado. A eficácia indica se o usuário consegue concluir a tarefa;
a eficiência, quanto tempo e esforço ele gasta para isso; e a
satisfação, como ele se sente ao usar o sistema.

Essas três dimensões mostram que um sistema não é fácil de usar apenas
por funcionar: ele
precisa ser compreendido logo no primeiro uso, evitar que o usuário
cometa erros e permitir que ele se recupere quando os erros acontecem.
No sistema de reserva de laboratórios, por exemplo, uma regra correta,
como o limite de 1 a 4 horas por reserva, pode gerar frustração se o
professor só descobrir esse limite depois de preencher todo o
formulário.

\subsection*{Teste de usabilidade com usuários}

Barbosa e Silva (2010) distinguem os métodos de avaliação feitos por
especialistas, que inspecionam a interface, dos métodos de observação,
em que pessoas representativas do público-alvo usam o sistema. No teste
de usabilidade, o avaliador prepara tarefas realistas, observa os
participantes enquanto eles as executam e registra dados como o sucesso
na tarefa, o tempo gasto e os erros cometidos. É comum pedir que o
participante pense em voz alta, explicando o que está tentando fazer, o
que ajuda a entender a causa das dificuldades, e não apenas o seu
efeito.

A quantidade de participantes é uma dúvida frequente. Nielsen e
Landauer (1993) propuseram um modelo matemático em que a proporção de
problemas encontrados depois de $i$ participantes é
$1-(1-\lambda)^{i}$, sendo $\lambda$ a probabilidade de um único
participante revelar um problema qualquer. Nos projetos analisados
pelos autores, $\lambda$ ficou, em média, em torno de 31\%. Com esse
valor, cinco participantes revelam cerca de 85\% dos problemas, e cada
participante a mais traz cada vez menos descobertas novas. Os autores
também mostram, numa análise de custo e benefício, que o melhor retorno
está em testes com poucos participantes, repetidos depois de cada
rodada de correções. O resultado, porém, depende de $\lambda$, que muda
de um sistema para outro; os cinco participantes são uma referência
prática, e não uma garantia.

\subsection*{Questionário SUS}

A \emph{System Usability Scale} (SUS) foi proposta por Brooke (1996)
como uma forma rápida e barata de medir a percepção de usabilidade de um
sistema. O questionário tem dez afirmações, respondidas numa escala de 1
(discordo totalmente) a 5 (concordo totalmente), alternando afirmações
positivas (itens ímpares) e negativas (itens pares). Para calcular a
nota, subtrai-se 1 da resposta de cada item ímpar e subtrai-se de 5 a
resposta de cada item par; a soma desses valores é multiplicada por
2,5, o que gera uma nota de 0 a 100 (BROOKE, 1996). Essa nota não é uma
porcentagem: ela só tem sentido quando comparada a uma referência.

Essa referência foi construída por Bangor, Kortum e Miller (2008), que
analisaram mais de dois mil questionários SUS aplicados em cerca de
duzentos testes de usabilidade, ao longo de dez anos. Os autores
confirmaram a alta confiabilidade do questionário e encontraram uma nota
média próxima de 70. Com base nisso, propuseram faixas de
aceitabilidade: notas abaixo de 50 indicam usabilidade inaceitável;
entre 50 e 70, marginal; e acima de 70, aceitável (BANGOR; KORTUM;
MILLER, 2008). Neste estudo, essas faixas são usadas para interpretar a
nota obtida pelo sistema.

\subsection*{Heurísticas de usabilidade}

Nielsen e Molich (1990) propuseram a avaliação heurística, em que
avaliadores examinam a interface com base em um pequeno conjunto de
princípios gerais, chamados heurísticas. As nove heurísticas originais
são: diálogo simples e natural; falar a língua do usuário; minimizar a
carga de memória do usuário; ser consistente; fornecer \emph{feedback};
oferecer saídas claramente marcadas; oferecer atalhos; fornecer boas
mensagens de erro; e prevenir erros. Os autores observaram que cada
avaliador, sozinho, encontrou apenas de 20\% a 51\% dos problemas, e
que juntar os resultados de vários avaliadores aumenta bastante a
quantidade de problemas encontrados (NIELSEN; MOLICH, 1990).

Neste trabalho, as heurísticas não são usadas como um método de
avaliação separado, mas como vocabulário para classificar os problemas
observados no teste com usuários. Cada problema também é classificado
pelo impacto que teve na tarefa observada: \emph{alto}, quando impediu o
participante de concluí-la; \emph{médio}, quando a atrasou ou exigiu uma
nova tentativa; e \emph{baixo}, quando causou apenas dúvida ou
incômodo. Essa classificação ajuda a decidir quais correções devem
entrar primeiro na versão 2.0.

\subsection*{Síntese}

Combinadas, essas referências orientam o método deste estudo: as tarefas
observadas medem a eficácia (taxa de sucesso) e a eficiência (tempo) do
uso; o questionário SUS mede a satisfação; e as heurísticas, com a
classificação por impacto, organizam os problemas encontrados numa
lista de melhorias para a versão 2.0.

\subsection*{Referências citadas na fundamentação}

{\small
BANGOR, A.; KORTUM, P. T.; MILLER, J. T. An empirical evaluation of the System Usability Scale. \emph{International Journal of Human-Computer Interaction}, v. 24, n. 6, p. 574--594, 2008. DOI: 10.1080/10447310802205776.

BARBOSA, S. D. J.; SILVA, B. S. \emph{Interação humano-computador}. Rio de Janeiro: Elsevier, 2010.

BROOKE, J. SUS: a ``quick and dirty'' usability scale. In: JORDAN, P. W. et al. (Ed.). \emph{Usability Evaluation in Industry}. London: Taylor \& Francis, 1996. p. 189--194.

NIELSEN, J.; LANDAUER, T. K. A mathematical model of the finding of usability problems. In: \emph{Proceedings of the INTERCHI'93 Conference on Human Factors in Computing Systems}. Amsterdam: ACM, 1993. p. 206--213. DOI: 10.1145/169059.169166.

NIELSEN, J.; MOLICH, R. Heuristic evaluation of user interfaces. In: \emph{Proceedings of the SIGCHI Conference on Human Factors in Computing Systems (CHI'90)}. Seattle: ACM, 1990. p. 249--256. DOI: 10.1145/97243.97281.
}

\end{document}
